\chapter{2011年大纲卷（文）}

\section{单选题}

\begin{problem}
设集合 \(U=\{1,2,3,4\}\)，\(M=\{1,2,3\}\)，\(N=\{2,3,4\}\)，则 \(\complement_U(M\cap N)=\)（\quad）

\choices
    {\(\{1,2\}\)}
    {\(\{2,3\}\)}
    {\(\{2,4\}\)}
    {\(\{1,4\}\)}
\end{problem}

\begin{answer}
D
\end{answer}

\begin{problem}
函数 \(y=2\sqrt{x}\)（\(x\geqslant0\)）的反函数为（\quad）

\choices
    {\(y=\frac{x^2}{4}\)（\(x\in\R\)）}
    {\(y=\frac{x^2}{4}\)（\(x\geqslant0\)）}
    {\(y=4x^2\)（\(x\in\R\)）}
    {\(y=4x^2\)（\(x\geqslant0\)）}
\end{problem}

\begin{answer}
B
\end{answer}

\begin{problem}
设向量 \(\bs{a}\)，\(\bs{b}\) 满足 \(|\bs{a}|=|\bs{b}|=1\)，\(\bs{a}\cdot\bs{b}=-\frac{1}{2}\)，则 \(|\bs{a}+2\bs{b}|=\)（\quad）

\choices
    {\(\sqrt{2}\)}
    {\(\sqrt{3}\)}
    {\(\sqrt{5}\)}
    {\(\sqrt{7}\)}
\end{problem}

\begin{answer}
B
\end{answer}

\begin{problem}
若变量 \(x\)，\(y\) 满足约束条件
\(\begin{cases}
x+y\leqslant6,\\
x-3y\leqslant-2,\\
x\geqslant1
\end{cases}\)，则 \(z=2x+3y\) 的最小值为（\quad）

\choices
    {\(17\)}
    {\(14\)}
    {\(5\)}
    {\(3\)}
\end{problem}

\begin{answer}
C
\end{answer}

\begin{problem}
下面四个条件中，使 \(a>b\) 成立的充分而不必要的条件是（\quad）

\choices
    {\(a>b+1\)}
    {\(a>b-1\)}
    {\(a^2>b^2\)}
    {\(a^3>b^3\)}
\end{problem}

\begin{answer}
A
\end{answer}

\begin{problem}
设 \(S_n\) 为等差数列 \(\{a_n\}\) 的前 \(n\) 项和. 若 \(a_1=1\)，公差 \(d=2\)，\(S_{k+2}-S_k=24\)，则 \(k=\)（\quad）

\choices
    {\(8\)}
    {\(7\)}
    {\(6\)}
    {\(5\)}
\end{problem}

\begin{answer}
D
\end{answer}

\begin{problem}
设函数 \(f(x)=\cos\omega x\)（\(\omega>0\)），将 \(y=f(x)\) 的图象向右平移 \(\frac{\pi}{3}\) 个单位长度后，所得的图象与原图象重合，则 \(\omega\) 的最小值等于（\quad）

\choices
    {\(\frac{1}{3}\)}
    {\(3\)}
    {\(6\)}
    {\(9\)}
\end{problem}

\begin{answer}
C
\end{answer}

\begin{problem}
已知直二面角 \(\alpha-l-\beta\)，点 \(A\in\alpha\)，\(AC\perp l\)，\(C\) 为垂足，点 \(B\in\beta\)，\(BD\perp l\)，\(D\) 为垂足. 若 \(AB=2\)，\(AC=BD=1\)，则 \(CD=\)（\quad）

\choices
    {\(2\)}
    {\(\sqrt{3}\)}
    {\(\sqrt{2}\)}
    {\(1\)}
\end{problem}

\begin{answer}
C
\end{answer}

\begin{problem}
\(4\) 位同学每人从甲，乙，丙 \(3\) 门课程中选修 \(1\) 门，则恰有 \(2\) 人选修课程甲的不同选法共有（\quad）

\choices
    {\(12\) 种}
    {\(24\) 种}
    {\(30\) 种}
    {\(36\) 种}
\end{problem}

\begin{answer}
B
\end{answer}

\begin{problem}
设 \(f(x)\) 是周期为 \(2\) 的奇函数，当 \(0\leqslant x\leqslant1\) 时，\(f(x)=2x(1-x)\)，则 \(f\left(-\frac{5}{2}\right)=\)（\quad）

\choices
    {\(-\frac{1}{2}\)}
    {\(-\frac{1}{4}\)}
    {\(\frac{1}{4}\)}
    {\(\frac{1}{2}\)}
\end{problem}

\begin{answer}
A
\end{answer}

\begin{problem}
设两圆 \(C_1\)，\(C_2\) 都和两坐标轴相切，且都过点 \((4,1)\)，则两圆心的距离 \(|C_1C_2|=\)（\quad）

\choices
    {\(4\)}
    {\(4\sqrt{2}\)}
    {\(8\)}
    {\(8\sqrt{2}\)}
\end{problem}

\begin{answer}
C
\end{answer}

\begin{problem}
已知平面 \(\alpha\) 截一球面得圆 \(M\)，过圆心 \(M\) 且与 \(\alpha\) 成 \(60^{\circ}\) 二面角的平面 \(\beta\) 截该球面得圆 \(N\). 若该球面的半径为 \(4\)，圆 \(M\) 的面积为 \(4\pi\)，则圆 \(N\) 的面积为（\quad）

\choices
    {\(7\pi\)}
    {\(9\pi\)}
    {\(11\pi\)}
    {\(13\pi\)}
\end{problem}

\begin{answer}
D
\end{answer}

\section{填空题}

\begin{problem}
\((1-x)^{10}\) 的二项展开式中，\(x\) 的系数与 \(x^9\) 的系数之差为\fillinblank{}.
\end{problem}

\begin{answer}
0
\end{answer}

\begin{problem}
已知 \(\alpha\in\left(\pi,\frac{3\pi}{2}\right)\)，\(\tan\alpha=2\)，则 \(\cos\alpha=\)\fillinblank{}.
\end{problem}

\begin{answer}
\(-\frac{\sqrt{5}}{5}\)
\end{answer}

\begin{problem}
已知正方体 \(ABCD-A_1B_1C_1D_1\) 中，\(E\) 为 \(C_1D_1\) 的中点，则异面直线 \(AE\) 与 \(BC\) 所成角的余弦值为\fillinblank{}.
\end{problem}

\begin{answer}
\(\frac{2}{3}\)
\end{answer}

\begin{problem}
已知 \(F_1\)，\(F_2\) 分别为双曲线 \(C:\frac{x^2}{9}-\frac{y^2}{27}=1\) 的左，右焦点，点 \(A\in C\)，点 \(M\) 的坐标为 \((2,0)\). \(AM\) 为 \(\angle F_1AF_2\) 的平分线. 则 \(|AF_2|=\)\fillinblank{}.
\end{problem}

\begin{answer}
6
\end{answer}

\section{解答题}

\begin{problem}
设等比数列 \(\{a_n\}\) 的前 \(n\) 项和为 \(S_n\)，已知 \(a_2=6\)，\(6a_1+a_3=30\)，求 \(a_n\) 和 \(S_n\).
\end{problem}

\begin{solution}
设等比数列的公比为 \(q\). 由 \(a_2=a_1q=6\)，且 \(a_3=a_2q=6q\)，已知条件给出 \(6a_1+6q=30\)，即 \(a_1+q=5\). 因此 \(a_1\) 与 \(q\) 是方程 \(t^2-5t+6=0\) 的两个根，所以 \(\{a_1,q\}=\{3,2\}\).

当 \(a_1=3\)，\(q=2\) 时，\(a_n=3\cdot2^{n-1}\)，从而 \(S_n=3(1+2+\cdots+2^{n-1})=3(2^n-1)\).

当 \(a_1=2\)，\(q=3\) 时，\(a_n=2\cdot3^{n-1}\)，从而 \(S_n=2(1+3+\cdots+3^{n-1})=3^n-1\). 综上，\(a_n=3\cdot2^{n-1}\)，\(S_n=3(2^n-1)\)，或 \(a_n=2\cdot3^{n-1}\)，\(S_n=3^n-1\).
\end{solution}

\begin{problem}
\(\triangle ABC\) 的内角 \(A\)，\(B\)，\(C\) 的对边分别为 \(a\)，\(b\)，\(c\). 已知 \(a\sin A+c\sin C-\sqrt{2}a\sin C=b\sin B\).

\begin{enumerate}
\item 求 \(B\)；
\item 若 \(A=75^{\circ}\)，\(b=2\)，求 \(a\)，\(c\).
\end{enumerate}
\end{problem}

\begin{solution}
\begin{enumerate}
\item 设三角形 \(ABC\) 的外接圆半径为 \(R\). 由正弦定理，\(a=2R\sin A\)，\(b=2R\sin B\)，\(c=2R\sin C\). 将其代入已知等式并约去 \(2R\)，得 \(\sin^2A+\sin^2C-\sqrt{2}\sin A\sin C=\sin^2B\). 又因为 \(A+B+C=\pi\)，所以 \(\sin B=\sin(A+C)\). 利用恒等式
\[
\sin^2A+\sin^2C-\sin^2(A+C)=-2\sin A\sin C\cos(A+C)
\]
可得
\[
0=\sin^2A+\sin^2C-\sqrt{2}\sin A\sin C-\sin^2(A+C)=\sin A\sin C\left[-2\cos(A+C)-\sqrt{2}\right]
\]
因为 \(A\)，\(C\) 是三角形内角，\(\sin A\sin C>0\)，故 \(\cos(A+C)=-\frac{\sqrt{2}}{2}\). 由于 \(A+C=\pi-B\)，有 \(-\cos B=-\frac{\sqrt{2}}{2}\)，且 \(0<B<\pi\)，所以 \(B=45^{\circ}\).

\item 当 \(A=75^{\circ}\) 时，\(C=180^{\circ}-75^{\circ}-45^{\circ}=60^{\circ}\). 再由正弦定理，\(a=\frac{b\sin A}{\sin B}=\frac{2\sin75^{\circ}}{\sin45^{\circ}}=1+\sqrt{3}\)，\(c=\frac{b\sin C}{\sin B}=\frac{2\sin60^{\circ}}{\sin45^{\circ}}=\sqrt{6}\).
\end{enumerate}
\end{solution}

\begin{problem}
根据以往统计资料，某地车主购买甲种保险的概率为 \(0.5\)，购买乙种保险但不购买甲种保险的概率为 \(0.3\). 设各车主购买保险相互独立.

\begin{enumerate}
\item 求该地 \(1\) 位车主至少购买甲，乙两种保险中的 \(1\) 种的概率；
\item 求该地 \(3\) 位车主中恰有 \(1\) 位车主甲，乙两种保险都不购买的概率.
\end{enumerate}
\end{problem}

\begin{solution}
\begin{enumerate}
\item 设事件 \(A\) 为车主购买甲种保险，事件 \(B\) 为车主购买乙种保险. 由于 \(B\cap\overline A\) 与 \(A\) 互不相容，且 \(B\cap\overline A\subset A\cup B\)，有 \(P(A\cup B)=P(A)+P(B\cap\overline A)=0.5+0.3=0.8\). 因此所求概率为 \(0.8\).
\item 一位车主两种保险都不购买的概率为 \(1-P(A\cup B)=0.2\). 由车主购买保险相互独立，恰有 \(1\) 位车主两种保险都不购买的概率为 \(\mathrm C_3^1(0.2)(1-0.2)^2=3\times0.2\times0.8^2=0.384\).
\end{enumerate}
\end{solution}

\begin{problem}
如图，四棱锥 \(S-ABCD\) 中，\(AB\parallel CD\)，\(BC\perp CD\)，侧面 \(SAB\) 为等边三角形. \(AB=BC=2\)，\(CD=SD=1\).

\begin{enumerate}
\item 证明：\(SD\perp\) 平面 \(SAB\)；
\item 求 \(AB\) 与平面 \(SBC\) 所成角的大小.
\end{enumerate}

\bitmapfigure[max width=0.55\linewidth]{img/2011/syllabus_liberal/q20_fig1.png}
\end{problem}

\begin{solution}
\begin{enumerate}
\item 建立直角坐标系，使底面 \(ABCD\) 在 \(z=0\) 平面内，取 \(B=(0,0,0)\)，\(A=(2,0,0)\)，\(C=(0,2,0)\)，\(D=(1,2,0)\). 设 \(S=(x,u,v)\). 因为三角形 \(SAB\) 为边长 \(2\) 的等边三角形，\(SA=SB=2\)，于是 \((x-2)^2+u^2+v^2=x^2+u^2+v^2=4\)，从而 \(x=1\)，\(u^2+v^2=3\). 又因 \(SD=1\)，有 \((u-2)^2+v^2=1\). 两式相减得 \(-4u+4=-2\)，所以 \(u=\frac{3}{2}\)，进而 \(v^2=\frac{3}{4}\). 取 \(v=\frac{\sqrt{3}}{2}\)，则
\(\overrightarrow{DS}=(0,-\frac{1}{2},\frac{\sqrt{3}}{2})\)，\(\overrightarrow{AS}=(-1,\frac{3}{2},\frac{\sqrt{3}}{2})\)，\(\overrightarrow{BS}=(1,\frac{3}{2},\frac{\sqrt{3}}{2})\). 因此 \(\overrightarrow{DS}\cdot\overrightarrow{AS}=0\times(-1)-\frac{1}{2}\times\frac{3}{2}+\frac{\sqrt{3}}{2}\times\frac{\sqrt{3}}{2}=0\)，\(\overrightarrow{DS}\cdot\overrightarrow{BS}=0\times1-\frac{1}{2}\times\frac{3}{2}+\frac{\sqrt{3}}{2}\times\frac{\sqrt{3}}{2}=0\)，即 \(SD\perp SA\)，\(SD\perp SB\). 由于 \(SA\)，\(SB\) 是平面 \(SAB\) 内相交的两条直线，所以 \(SD\perp\) 平面 \(SAB\).

\item 取平面 \(SBC\) 的一个法向量 \(\bs{n}=\overrightarrow{BS}\times\overrightarrow{BC}=(-\sqrt{3},0,2)\)，并取直线 \(AB\) 的方向向量 \(\overrightarrow{BA}=(2,0,0)\). 设 \(AB\) 与平面 \(SBC\) 所成的角为 \(\theta\)，则
\[
\sin\theta=\frac{|\overrightarrow{BA}\cdot\bs n|}{|\overrightarrow{BA}|\,|\bs n|}=\frac{2\sqrt{3}}{2\sqrt{7}}=\frac{\sqrt{21}}{7}
\]
所以 \(\theta=\arcsin\frac{\sqrt{21}}{7}\).
\end{enumerate}
\end{solution}

\begin{problem}
已知函数 \(f(x)=x^3+3ax^2+(3-6a)x+12a-4\)（\(a\in\R\)）.

\begin{enumerate}
\item 证明：曲线 \(y=f(x)\) 在 \(x=0\) 处的切线过点 \((2,2)\)；
\item 若 \(f(x)\) 在 \(x=x_0\) 处取得极小值，\(x_0\in(1,3)\)，求 \(a\) 的取值范围.
\end{enumerate}
\end{problem}

\begin{solution}
\begin{enumerate}
\item 先求导得 \(f'(x)=3x^2+6ax+3-6a\)，所以 \(f(0)=12a-4\)，\(f'(0)=3-6a\). 曲线在 \(x=0\) 处的切线方程为 \(y=f(0)+f'(0)x=(12a-4)+(3-6a)x\). 当 \(x=2\) 时，\(y=12a-4+2(3-6a)=2\)，故该切线过点 \((2,2)\).
\item 先求导得 \(f'(x)=3x^2+6ax+3-6a=3\bigl(x^2+2ax+1-2a\bigr)\). 曲线在 \(x=x_0\) 处取得极小值时，必须有 \(f'(x_0)=0\). 由于 \(x_0\in(1,3)\)，由
\(x_0^2+2ax_0+1-2a=0\) 解得 \(a=-\frac{x_0^2+1}{2(x_0-1)}\).

二次式 \(f'(x)/3\) 的二次项系数为正，极小值对应较大的临界点，等价于 \(f''(x_0)=6(x_0+a)>0\). 代入上式，得
\[
x_0+a=x_0-\frac{x_0^2+1}{2(x_0-1)}=\frac{x_0^2-2x_0-1}{2(x_0-1)}>0
\]
由于 \(x_0>1\)，所以 \(x_0>1+\sqrt{2}\). 因而 \(x_0\in(1+\sqrt{2},3)\).

令 \(a(x)=-\frac{x^2+1}{2(x-1)}\)，则在 \(x>1+\sqrt{2}\) 时
   \(a'(x)=-\frac{x^2-2x-1}{2(x-1)^2}<0\). 所以当 \(x\) 从 \(1+\sqrt{2}\) 增大到 \(3\) 时，\(a(x)\) 严格减小，且 \(a(1+\sqrt{2})=-1-\sqrt{2}\)，\(a(3)=-\frac{5}{2}\). 由于两个端点均不取，且该函数在区间内连续，得到 \(a\in\left(-\frac{5}{2},-1-\sqrt{2}\right)\). 反之，区间内每个 \(a\) 都对应唯一的 \(x_0\in(1+\sqrt{2},3)\)，并满足 \(f''(x_0)>0\)，故该范围即为所求.
\end{enumerate}
\end{solution}

\section{选考题：解答题}

\begin{problem}
已知 \(O\) 为坐标原点，\(F\) 为椭圆 \(C:x^2+\frac{y^2}{2}=1\) 在 \(y\) 轴正半轴上的焦点，过 \(F\) 且斜率为 \(-\sqrt{2}\) 的直线 \(l\) 与 \(C\) 交于 \(A\)，\(B\) 两点，点 \(P\) 满足 \(\overrightarrow{OA}+\overrightarrow{OB}+\overrightarrow{OP}=0\).

\begin{enumerate}
\item 证明：点 \(P\) 在 \(C\) 上；
\item 设点 \(P\) 关于点 \(O\) 的对称点为 \(Q\)，证明：\(A\)，\(P\)，\(B\)，\(Q\) 四点在同一圆上.
\end{enumerate}

\FigureLayoutDeclare{img/2011/syllabus_liberal/q22_fig1.png}{6.0cm}{16.801bp 16.801bp 16.801bp 16.801bp}
\bitmapfigure[max width=0.38\linewidth]{img/2011/syllabus_liberal/q22_fig1.png}
\end{problem}

\begin{solution}
\begin{enumerate}
\item 椭圆方程可写为 \(\frac{x^2}{1^2}+\frac{y^2}{(\sqrt{2})^2}=1\)，故其焦半距为 \(c=\sqrt{(\sqrt{2})^2-1^2}=1\)，且焦点 \(F=(0,1)\). 直线 \(l\) 过 \(F\)，斜率为 \(-\sqrt{2}\)，所以方程为 \(y=1-\sqrt{2}x\). 将其代入椭圆方程，得
\[
x^2+\frac{(1-\sqrt{2}x)^2}{2}=1
\]
即 \(4x^2-2\sqrt{2}x-1=0\). 按图中位置，\(A\) 的横坐标取较小根，\(B\) 的横坐标取较大根，因此
\(x_A=\frac{\sqrt{2}-\sqrt{6}}{4}\)，\(x_B=\frac{\sqrt{2}+\sqrt{6}}{4}\)，由直线方程得 \(y_A=\frac{1+\sqrt{3}}{2}\)，\(y_B=\frac{1-\sqrt{3}}{2}\). 由 \(\overrightarrow{OA}+\overrightarrow{OB}+\overrightarrow{OP}=0\)，得到 \(P=-(A+B)=(-\frac{\sqrt{2}}{2},-1)\). 于是
\(x_P^2+\frac{y_P^2}{2}=\frac{1}{2}+\frac{1}{2}=1\)，所以点 \(P\) 在椭圆 \(C\) 上.
\item 设 \(r=\frac{\sqrt{2}}{2}\)，则 \(P=(-r,-1)\)，点 \(P\) 关于原点 \(O\) 的对称点为 \(Q=(r,1)\). 考虑圆
\[
\Gamma:x^2+y^2+\frac{1}{2\sqrt{2}}x-\frac{1}{4}y-\frac{3}{2}=0.
\]
将 \(P\)，\(Q\) 的坐标代入，利用 \(r^2=\frac12\)，分别得到 \(\frac32-\frac{r^2}{2}+\frac14-\frac32=0\) 与 \(\frac32+\frac{r^2}{2}-\frac14-\frac32=0\)，所以两点均在圆 \(\Gamma\) 上. 对于 \(A\)，\(B\)，由直线方程 \(y=1-\sqrt{2}x=1-2rx\) 以及椭圆方程，有
\[
x^2+y^2+\frac{r}{2}x-\frac14y-\frac32=\frac14-x^2+rx
\]
而交点横坐标满足 \(x^2-rx-\frac14=0\)，故上式等于 \(0\). 因此 \(A\)，\(B\) 也都在圆 \(\Gamma\) 上，故 \(A\)，\(P\)，\(B\)，\(Q\) 四点在同一圆上.
\end{enumerate}
\end{solution}
